% =====================================================================
%  CHAPTER 4 — SIMULATION RESULTS AND ANALYSIS
% =====================================================================
\chapter{SIMULATION RESULTS AND ANALYSIS}
\label{ch:results}

The previous chapters built the simulation from three foundations: the physical
theory of solar radio interferometry in Chapter~\ref{ch:theory}, the receiving-element
model in Chapter~\ref{ch:antenna}, and the end-to-end software pipeline in
Chapter~\ref{ch:pipeline}. This chapter evaluates whether those components work
together as a useful design tool for a future low-cost solar radio interferometer.

The purpose of this chapter is not to introduce new theory or new derivations; those
belong to Chapters~\ref{ch:theory}, \ref{ch:antenna} and~\ref{ch:pipeline}. Instead,
each result below uses the relations already established there as a validation
benchmark and answers one practical question about the simulation pipeline: whether
the geometry, the source model, the imaging stage, the array layout, the receiver
sensitivity, and the antenna beam are all correctly and usefully represented.
Agreement with an analytic benchmark does not mean that the Learmonth receiver
hardware itself has been reproduced; it means that the proposed low-cost
interferometer model responds correctly when driven by a real solar flux record.
The results are therefore presented as a sequence of validation levels, moving from
analytic checks (Sections~\ref{sec:fringe-results}-\ref{sec:dirty-results}) to design
and sensitivity studies (Sections~\ref{sec:sweep}-\ref{sec:beam-fidelity}), and
finally to practical deployment diagnostics (Sections~\ref{sec:apsynsim}-\ref{sec:clock-offset}).

% ---------------------------------------------------------------------
\section{Result Strategy and Validation Logic}
\label{sec:result-strategy}

Table~\ref{tab:result-strategy} summarises the logic followed throughout the chapter:
each row identifies a block of the simulation pipeline, the result used to test it,
and what agreement with that result actually establishes. Reading the chapter with
this table in mind clarifies why the results are ordered as they are - each section
depends on the blocks validated by the sections before it.

\begin{table}[H]
  \centering
  \caption[Validation strategy: pipeline block, test result, and what it establishes]{Validation strategy: pipeline block, test result, and what it establishes.\\[4pt]
  Each result in this chapter targets one block of the simulation pipeline
  described in Chapter~\ref{ch:pipeline}. Sections~\ref{sec:fringe-results}-\ref{sec:dirty-results}
  test internal correctness against analytic predictions; Sections~\ref{sec:sweep}-\ref{sec:beam-fidelity}
  use the validated pipeline for design and sensitivity studies; Sections~\ref{sec:apsynsim}-\ref{sec:clock-offset}
  place the simulation in context and diagnose the gap to real hardware.}
  \label{tab:result-strategy}
  \renewcommand{\arraystretch}{1.3}
  \begin{tabularx}{\textwidth}{l X X}
    \toprule
    \textbf{Pipeline block} & \textbf{Result used to test it} & \textbf{What it establishes} \\
    \midrule
    Geometry (ENU$\rightarrow$XYZ$\rightarrow\uvw$, hour angle)
      & Fringe period, Section~\ref{sec:fringe-results}
      & Baseline and hour-angle geometry are implemented correctly \\
    Source model (resolved solar disk)
      & $\jinc$ visibility decay, Section~\ref{sec:vis-results}
      & Solar angular size is represented correctly \\
    Correlator, $uv$-sampling, inverse DFT
      & Dirty image $=$ sky $*$ PSF, Section~\ref{sec:dirty-results}
      & The imaging stage is internally consistent \\
    Array layout
      & $uv$-coverage, PSF, dynamic range across five geometries, Section~\ref{sec:sweep}
      & Which layout is preferable for a future deployment \\
    Receiver noise model
      & Minimum detectable flux and burst SNR, Section~\ref{sec:snr-analysis}
      & Quiet Sun and bursts are detectable under the model's assumptions \\
    Antenna primary beam
      & Gaussian versus 4NEC2 transit power, Section~\ref{sec:beam-fidelity}
      & The real antenna beam materially changes calibration, not just cosmetics \\
    Field deployment
      & Solar-excess comparison with Hoa~Lac hardware, Section~\ref{sec:real-comparison}
      & Where the gap between simulation and hardware measurement originates \\
    Clock model
      & Phase-error and coherence-loss study, Section~\ref{sec:clock-offset}
      & Why unsynchronised SDRs cannot recover a stable fringe \\
    \bottomrule
  \end{tabularx}
\end{table}

The remainder of this chapter follows this order: each section states the question
being tested, presents the result, interprets it against the corresponding block of
the pipeline, and, where relevant, draws out its implication for a future hardware
deployment.

% ---------------------------------------------------------------------
\section{Geometry Validation: Fringe Pattern}
\label{sec:fringe-results}

The first and most direct question the pipeline must answer is whether the array
geometry and hour-angle calculation are correct. Driving the two-element sub-array at
the \SI{8}{\metre} baseline with the real flux and the exact $\uvw$ geometry
produces the cross-correlation amplitude shown in Fig.~\ref{fig:fringe}. The
measured fringe period is approximately 14~minutes, in agreement with the analytic
prediction of Eq.~\eqref{eq:fringeperiod} for $b=\SI{8}{\metre}$,
$\lambda=\SI{0.49}{\metre}$ and $\delta=0.0^\circ$ (spring equinox). Because the fringe period depends
only on the geometry and not on the source structure, this agreement confirms that the
ENU$\rightarrow$XYZ$\rightarrow\uvw$ transformation~\citep{tms2017} and the hour-angle parameterisation
are implemented correctly, and that the pipeline can be trusted to test baseline and
array-layout choices ahead of a physical deployment.

\begin{figure}[H]
  \centering
  \includegraphics[width=\textwidth]{fig_fringe_pattern}
  \caption[Simulated fringe pattern for the 8-m baseline (20~March~2025)]{Simulated fringe pattern for the 8-m baseline (20~March~2025).\\[4pt]
  The real part of the baseline visibility is plotted over the full
  spring-equinox transit ($H\in[\pm90^\circ]$, $\delta=0^\circ$). The
  approximately 14-minute fringe period matches the analytic prediction of
  Eq.~\eqref{eq:fringeperiod} to within the sampling resolution, confirming
  that the ENU$\rightarrow$XYZ$\rightarrow\uvw$ transformation and the
  hour-angle parameterisation are both correctly implemented.}
  \label{fig:fringe}
\end{figure}

% ---------------------------------------------------------------------
\section{Source-Structure Validation: The Resolved Solar Disk}
\label{sec:vis-results}

The Sun is not a point source, so the next question is whether the simulated array
correctly represents its finite angular size. When the source is modelled as a
uniformly bright disk of the solar angular size, the
amplitude of the visibility as a function of baseline length follows the $\jinc$ curve of
Eq.~\eqref{eq:jinc}. Figure~\ref{fig:jinc} overlays the simulated visibility amplitudes
on the analytic $\jinc$ function; the simulated points fall on the curve, and the first
null occurs at the predicted baseline. This confirms that the forward visibility engine
correctly implements the van Cittert-Zernike relation~\citep{tms2017,taylor1999synthesis} and that the array, when extended,
is able to resolve the solar disk rather than merely detect it.

\begin{figure}[H]
  \centering
  \includegraphics[width=\textwidth]{fig_visibility_jinc}
  \caption[Visibility amplitude versus baseline length: simulation versus $\jinc$ prediction]{Visibility amplitude versus baseline length: simulation versus $\jinc$ prediction.\\[4pt]
  Simulated visibility amplitudes (points) at the transit peak are overlaid on
  the analytic $\jinc(u\theta_\odot)$ curve for a uniform solar disk of angular
  diameter approximately 32~arcminutes (Eq.~\eqref{eq:jinc}). Agreement across
  all sampled baselines, including the position of the first null, validates
  that the forward visibility engine correctly implements the van
  Cittert-Zernike theorem.}
  \label{fig:jinc}
\end{figure}

% ---------------------------------------------------------------------
\section{Imaging Validation: Dirty Image as Sky Convolved with the PSF}
\label{sec:dirty-results}

With the geometry and the source model both validated, the remaining question is
whether the imaging stage correctly turns sparse visibility samples into a
synthesised image. The central consistency check is Eq.~\eqref{eq:dirty}: the dirty
image must equal the input sky convolved with the independently computed dirty beam~\citep{taylor1999synthesis,tms2017}.
Figure~\ref{fig:psf-check} shows, side by side, the dirty image produced by the full
DFT inversion and the sky model convolved with the PSF computed from the sampling
function alone. The two are indistinguishable to numerical precision, demonstrating that
the correlator, the $uv$-sampling, and the inverse transform are mutually consistent.

\begin{figure}[H]
  \centering
  \begin{subfigure}{0.48\textwidth}
    \includegraphics[width=\textwidth]{fig_dirty_from_dft}
    \caption{Dirty image (DFT inversion)}
  \end{subfigure}\hfill
  \begin{subfigure}{0.48\textwidth}
    \includegraphics[width=\textwidth]{fig_sky_conv_psf}
    \caption{Sky $*$ PSF (independent)}
  \end{subfigure}
  \caption[Dirty-image consistency check: DFT inversion versus sky convolved with PSF]{Dirty-image consistency check: DFT inversion versus sky convolved with PSF.\\[4pt]
  Panel~(a) shows the dirty image produced by inverting all simulated
  visibilities via the direct DFT of Eq.~\eqref{eq:dirty-dft}. Panel~(b)
  shows the same sky model convolved with the independently computed
  synthesised beam. The two are numerically identical to machine precision,
  demonstrating that the correlator, the $uv$-sampling function, and the
  inverse DFT together implement Eq.~\eqref{eq:dirty} correctly.}
  \label{fig:psf-check}
\end{figure}

Two further points support the credibility of this three-level analytic validation
(geometry, source structure, imaging). First, the quiet-Sun \SI{610}{\mega\hertz} flux
used to drive the pipeline is stable over the transit; far from being a weakness, this
stability is ideal for instrument validation, since any departure of the simulated
fringe from the prediction can then be attributed to the modelled geometry rather than
to source variability. Second, the ionosphere is effectively common-mode for the
baselines used ($\lesssim\SI{130}{\metre}$ at \SI{610}{\mega\hertz}), with a residual
differential phase of about one degree, so it does not corrupt the differential
measurement on which the fringe and visibility checks depend. Taken together with
Sections~\ref{sec:fringe-results} and~\ref{sec:vis-results}, the credibility of the
simulation therefore rests not on the number of observing days but on the simultaneous
agreement with several independent analytic benchmarks - exactly the standard by which
an instrument model should be judged, and the basis on which the design and
sensitivity studies of the following sections can be trusted.

% ---------------------------------------------------------------------
\section{Array-Geometry Design Study}
\label{sec:sweep}

With the pipeline validated at the analytic level, it can be used to answer a design
question that cannot be tested cheaply in the field: which antenna layout is
preferable for a future deployment? The pipeline was run over a sweep of
five array geometries sharing the same number of elements: a compact east-west line, a
north-south line, a tee, a graded cross, and a randomised layout. For each, the
$uv$-coverage, the dirty beam, and the restored image were computed and exported as
static figures together with the numerical data. Table~\ref{tab:sweep} summarises the
qualitative trade-offs.

\begin{table}[H]
  \centering
  \caption[Qualitative comparison of array geometries in the parametric sweep]{Qualitative comparison of array geometries in the parametric sweep.\\[4pt]
  Same element count in all configurations). Each row summarises the dominant
  imaging characteristic of the corresponding layout: the compact east-west
  array matches the deployed field configuration but resolves only one axis;
  the graded cross achieves the highest two-dimensional coverage and the
  lowest PSF sidelobe level.}
  \label{tab:sweep}
  \renewcommand{\arraystretch}{1.25}
  \begin{tabularx}{\textwidth}{l l X}
    \toprule
    \textbf{Layout} & \textbf{Max baseline} & \textbf{Imaging characteristic} \\
    \midrule
    Compact EW    & \SI{24}{\metre}  & Matches the field array; resolves only the EW
                                       direction; elongated beam. \\
    NS line       & \SI{105}{\metre} & Complementary NS resolution; poor EW. \\
    Tee           & moderate         & Two-dimensional coverage; asymmetric beam. \\
    Graded cross  & \SI{130}{\metre} & Most isotropic coverage; resolves the solar disk;
                                       lowest sidelobes. \\
    Random        & variable         & Irregular but space-filling coverage; useful as a
                                       reference. \\
    \bottomrule
  \end{tabularx}
\end{table}

\begin{figure}[H]
  \centering
  \includegraphics[width=\textwidth]{fig_layout_sweep}
  \caption[$uv$-coverage and dirty beams for the five-geometry parametric sweep]{$uv$-coverage and dirty beams for the five-geometry parametric sweep.\\[4pt]
  Top row: $uv$-plane track patterns accumulated over the 12-hour transit;
  bottom row: corresponding synthesised dirty beams. Columns from left to
  right: compact east-west, north-south linear, tee-shaped, extended graded
  cross, and random two-dimensional. The graded cross (baselines to
  \SI{130}{\metre}) produces the most isotropic $uv$-coverage and the most
  compact, symmetric beam with the lowest sidelobe level.}
  \label{fig:sweep}
\end{figure}

The compact east-west array-the geometry actually built in the field-detects the Sun
and reproduces the correct fringe behaviour but, with baselines only to \SI{24}{\metre},
resolves structure in one direction only. The extended graded cross, reaching
\SI{130}{\metre}, achieves the two-dimensional coverage needed to resolve the solar disk.
This quantifies the upgrade path from the existing instrument to one capable of imaging,
rather than merely detecting, solar structure. Longer baselines alone are not
sufficient, however: PSF sidelobe level and $uv$-coverage uniformity are equally
decisive, as the dynamic-range comparison of Section~\ref{sec:detectability} shows. The
choice of a deployment layout must therefore weigh baseline length against these
attributes together, rather than maximising baseline length in isolation.

\paragraph{A minimum-redundancy alternative: Case~1.75.} Section~\ref{subsec:case175}
of Chapter~\ref{ch:pipeline} evaluated a sixth layout at the same element count and
spacing scale as Cases~1 and~1.5, but built on a different design principle: a
DLITE-inspired triangle-plus-centre arrangement~\citep{helmboldt2021dlite} in which
every one of the six baselines points in a distinct direction, rather than the two
redundant pairs produced by a symmetric cross. The result (Figs.~\ref{fig:case175-geometry}
and~\ref{fig:case175-eval} of Chapter~\ref{ch:pipeline}) confirms that this
non-redundant $uv$-sampling strategy is genuinely effective at extracting
two-dimensional image information from a minimal antenna count: the North-South FWHM
improves $9.1\times$ over the purely East-West Case~1, exceeding the resolution
recovered by the symmetric cross of Case~1.5 with the same four elements and the same
$8$-m spacing unit.

This benefit is intrinsically tied to using at least four elements: three antennas
alone can form only a single triangle with no independent centre-referenced baseline,
and fewer than four elements cannot supply the six mutually non-redundant $uv$-tracks
that make the layout effective in the first place - consistent with DLITE itself being
built from exactly four elements for this reason, no more and no fewer. It is equally
important not to over-read the absolute resolution achieved here, or to read Case~1.75
as a simulation of DLITE itself: Case~1.75 reuses the $8$-m baseline scale of Cases~1
and~1.5 purely for direct comparability, and is evaluated at \SI{610}{\mega\hertz},
whereas DLITE operates at \SIrange{30}{40}{\mega\hertz} with LWA-style crossed-dipole
antennas on baselines of several hundred metres~\citep{helmboldt2021dlite,carson2022dlite}.
Only the minimum-redundancy \emph{topology} is carried over here, not DLITE's
frequency, antenna, or baseline scale. The result therefore validates that topology
-non-redundant triangle-plus-centre sampling- as an efficient way to reconstruct a
two-dimensional image from four elements, not the absolute angular resolution that a
future physical deployment following the same principle could achieve with longer
baselines.

% ---------------------------------------------------------------------
\section{Sensitivity and Burst Detectability}
\label{sec:snr-analysis}

The design study of Section~\ref{sec:sweep} assumed that the receiver has enough
sensitivity to detect the source in the first place. This section answers that
question directly: is the solar signal strong enough, given the simulated receiver's
noise floor, to be detected and, in the case of transient bursts, resolved with
useful dynamic range?

\subsection{The Radiometer Equation}

The fundamental sensitivity limit of a single-dish (or single-baseline) radio
measurement is the \emph{radiometer equation}~\citep{tms2017,kraus1986}: for a
total-power integration over time $\tau$ and bandwidth $\Delta\nu$, the RMS
fluctuation of the system temperature estimate, $\Delta T_\mathrm{rms}$
(Eq.~\eqref{eq:radiometer} of Chapter~\ref{ch:theory}), converts via the Friis
relation (Eq.~\eqref{eq:friis}) to a minimum detectable flux density (MDF):
\begin{equation}
  \Delta S_{\rm min} = \frac{2\,k_B\,T_{\rm sys}}{A_e\sqrt{\Delta\nu\,\tau}}.
  \label{eq:mdf}
\end{equation}
For the simulation parameters ($T_{\rm sys}=\SI{150}{\kelvin}$,
$\Delta\nu=\SI{1}{\mega\hertz}$, $A_e\approx\SI{0.17}{\metre\squared}$),
Table~\ref{tab:snr-bins} evaluates Eq.~\eqref{eq:mdf} over a range of
\texttt{INTEGRATION\_BINS} values and compares $\Delta S_{\rm min}$ with
the quiet-Sun flux of $\sim\SI{100}{\sfu}$.

\begin{table}[H]
  \centering
  \caption[Minimum detectable flux density and SNR versus integration-dump count]{Minimum detectable flux density and SNR versus integration-dump count.\\[4pt]
  Values of $\Delta S_{\rm min}$ (Eq.~\eqref{eq:mdf}) and the quiet-Sun SNR
  are tabulated as a function of the number of correlator dumps
  (\texttt{INTEGRATION\_BINS})
  over the 12-hour spring-equinox transit.
  The configuration used in this thesis is highlighted.}
  \label{tab:snr-bins}
  \renewcommand{\arraystretch}{1.3}
  \begin{tabular}{c c c c c}
    \toprule
    \textbf{Bins} & $\boldsymbol{\tau}$ (s) & $\boldsymbol{\sigma_{\Delta T}}$ (mK)
      & $\boldsymbol{\Delta S_{\rm min}}$ (SFU) & \textbf{Quiet-Sun SNR} \\
    \midrule
    10   & 4320 & 2.3  & 0.026 & $\sim$3\,800 \\
    25   & 1728 & 3.6  & 0.041 & $\sim$2\,400 \\
    \rowcolor{LightCyan}
    50   & 864  & 5.1  & 0.058 & $\sim$1\,700 \\
    100  & 432  & 7.2  & 0.082 & $\sim$1\,200 \\
    200  & 216  & 10.2 & 0.116 & $\sim$860   \\
    500  & 86   & 16.2 & 0.183 & $\sim$540   \\
    \bottomrule
  \end{tabular}
\end{table}

All entries in Table~\ref{tab:snr-bins} yield quiet-Sun SNR in excess of
500:1. Even the coarsest choice (10 dumps) achieves
$\Delta S_{\rm min}\approx\SI{0.026}{\sfu}$, whereas the quiet-Sun flux
is $\sim\SI{100}{\sfu}$-a margin of more than $3\,700\times$. The
simulation is therefore not sensitivity-limited by thermal noise; it is
limited by the sparsity of the UV coverage and the associated PSF sidelobe
level, consistent with the layout trade-offs identified in Section~\ref{sec:sweep}.

\subsection{Per-Sample Signal-to-Noise Ratio}

At the level of a single voltage sample, the picture looks different. For the
quiet-Sun case ($S_\nu\approx\SI{100}{\sfu}$), the signal power at a single antenna
element follows from Eq.~\eqref{eq:psig}:
\begin{equation}
  P_{\rm sig} = \tfrac{1}{2}\times0.17\times10^{-22}\times100\times10^6
    \approx \SI{8.5e-16}{\watt}.
  \label{eq:psig-value}
\end{equation}
The corresponding Johnson-Nyquist noise power (Eq.~\eqref{eq:johnson}) is
\begin{equation}
  P_n = 1.381\times10^{-23}\times150\times10^6
    \approx \SI{2.07e-15}{\watt}.
  \label{eq:johnson-value}
\end{equation}
The per-sample SNR is thus
\begin{equation}
  \frac{P_{\rm sig}}{P_n} \approx 0.41 \;\approx\; -3.8~\si{dB},
  \label{eq:persample-snr}
\end{equation}
meaning each individual voltage sample is noise-dominated, as in any real radio
telescope. Averaging over $N_s$ samples per dump recovers the signal,
\begin{equation}
  \mathrm{SNR}_{\rm dump} = \frac{P_{\rm sig}}{P_n}\sqrt{N_s}
    \approx 0.41\times\sqrt{1000} \approx 13,
  \label{eq:snr-dump}
\end{equation}
where $N_s = N_{\rm samples}/\texttt{INTEGRATION\_BINS} = 50\,000/50 = 1\,000$ is the
number of simulated voltage samples averaged per correlator dump - a much smaller
count than the Nyquist-independent sample number $\Delta\nu\,\tau$ implicit in the
radiometer equation (Eq.~\eqref{eq:radiometer}) behind Table~\ref{tab:snr-bins}. This
gives a time-averaged fringe amplitude recoverable at $13\times\sqrt{50}\approx92$:1
across the full transit. This per-sample-count check and the radiometer-equation
sensitivity of Table~\ref{tab:snr-bins} answer two different questions - the former
confirms that the implemented discrete voltage chain behaves as expected at the
sample level, the latter states the physical sensitivity limit used throughout the
rest of this chapter - and the two figures should not be read as numerically
interchangeable. This low per-sample SNR is also why the sliding-RMS envelope
visualisation of Chapter~\ref{ch:pipeline} is necessary: only a $\sim\!500$-sample
running mean reveals the slowly varying solar-flux envelope hidden in a single
noise-like voltage trace.

\subsection{The LSB $\geq 3$ Criterion as a Quantisation Sanity Check}

Before FPGA export, the pipeline enforces the LSB check of
Eq.~\eqref{eq:lsb-check} from Chapter~\ref{ch:pipeline}: the noise voltage
$V_{\rm noise}$ multiplied by the ADC scale factor of Eq.~\eqref{eq:scale} must be
at least 3. If this fails, the noise is clipped
into fewer than three LSBs, meaning the ADC is over-ranging and quantisation
noise would corrupt the cross-correlation. For the March 2025 transit the noise
floor occupies well over 20 LSBs at the chosen scale factor, confirming that the
16-bit dynamic range is correctly matched to the signal-plus-noise envelope.

\subsection{Flux Levels and Detectability of Solar Radio Events}
\label{sec:detectability}

Solar radio emission at \SI{610}{\mega\hertz} spans more than four orders of
magnitude in flux density. Table~\ref{tab:detectability} classifies the principal
solar radio event types by their typical flux levels and evaluates the
detection significance against the simulation's sensitivity floor
($\Delta S_{\rm min}\approx\SI{0.058}{\sfu}$ for \texttt{INTEGRATION\_BINS}$=50$,
from Table~\ref{tab:snr-bins}). The quiet-Sun and burst flux values are
order-of-magnitude reference figures, not measurements specific to 20~March~2025:
they are drawn from the classical solar radio burst classification
literature~\citep{wild1963,dulk1985,mclean1985} and from directly reported
610~MHz values in~\citet{marque2018}, who cite a NOAA-reported peak flux of
$\sim\SI{1000}{\sfu}$ at \SI{610}{\mega\hertz} for a moderate event (the basis
for the ``Type~III (moderate)'' row) and a quiet-Sun flux of $\SI{75}{\sfu}$ at
the nearby decimetric frequency of \SI{1060}{\mega\hertz}. The quiet-Sun entry
of $\sim\SI{100}{\sfu}$ used throughout this chapter is accordingly a round,
slightly conservative reference value rather than a precise measurement: the
actual 20~March~2025 record (Fig.~\ref{fig:sws-flux} of Chapter~\ref{ch:pipeline})
shows a quieter baseline of approximately $\SIrange{65}{70}{\sfu}$, consistent with
this literature anchor. Because $\Delta S_{\rm min}$ is more than three orders of
magnitude below either figure, this choice of reference value has no effect on
the detectability conclusions below.

\begin{table}[H]
  \centering
  \caption[Solar event detectability at 610~MHz versus simulation noise floor]{Solar event detectability at 610~MHz versus simulation noise floor.\\[4pt]
  Representative solar radio event types are classified by their typical flux
  density at \SI{610}{\mega\hertz} and compared against the simulation's
  minimum detectable flux ($\Delta S_{\rm min}\approx0.058$~SFU for 50 dumps).
  Flux levels are order-of-magnitude reference values from
  \citet{wild1963,dulk1985,mclean1985,marque2018}, not measurements from the
  specific transit simulated in this thesis (see text). All event categories,
  including the quiet Sun at $\approx100$~SFU, are detectable at extremely high
  significance in a single integration dump. Type~III and Type~II bursts exceed
  the noise floor by factors of $10^4$--$10^5$.}
  \label{tab:detectability}
  \renewcommand{\arraystretch}{1.3}
  \begin{tabularx}{\textwidth}{l c c >{\raggedright\arraybackslash}X}
    \toprule
    \textbf{Event type} & \textbf{Typical flux (SFU)} & \textbf{SNR (1 dump)} & \textbf{Detectability} \\
    \midrule
    Quiet Sun               & $\sim$100           & $\sim$1\,700  & Always detected \\
    Low-level active region & $\sim$300           & $\sim$5\,200  & Easily detected \\
    Type~III (moderate)     & $\sim$1\,000        & $\sim$17\,200 & Strongly detected \\
    Type~III (intense)      & $\sim$10\,000       & $\sim$172\,000& Scale-saturating \\
    Type~II (shock)         & 500-5\,000         & 8\,600-86\,000 & Time-extended \\
    Type~IV (continuum)     & $\sim$200-1\,000   & 3\,400-17\,200 & Detected with high confidence \\
    \bottomrule
  \end{tabularx}
\end{table}

All event categories in Table~\ref{tab:detectability} are detectable at extremely
high significance in a single integration dump. The quiet Sun itself-the weakest
source in the table-registers at SNR $\approx1\,700$ per dump. This confirms that,
under the simulated receiver assumptions and when driven by the Learmonth
\SI{610}{\mega\hertz} flux record, quiet-Sun and burst-level fluxes are well above
the model noise floor, and that any Type~II or Type~III burst occurring during the
transit would be unmistakable under the model's assumptions.

\subsection{Dynamic Range of the Dirty Image}

The dynamic range (DR) of a dirty image is the ratio of the peak brightness to the
RMS of the off-source background, using the same definition as the array-geometry
sweep of Chapter~\ref{ch:pipeline}:
\begin{equation}
  \mathrm{DR} = \frac{I^D_{\rm peak}}{\sigma_{I^D}}.
  \label{eq:dynamic-range}
\end{equation}
For the compact east-west array (Case~1), the dominant image artefact is not
thermal noise but the north-south elongation of the dirty beam: the PSF peak
sidelobe is a significant fraction of the main-lobe amplitude, and it appears
as a systematic stripe rather than as random noise. Consequently, the image DR
is limited by the PSF sidelobe level rather than by receiver noise, echoing the
layout-versus-sensitivity trade-off already noted in Section~\ref{sec:sweep}.
Additional integration time does not help in this regime; only more baselines
(better UV coverage) or deconvolution (\textsc{clean}) can improve it.
Transitioning to the graded-cross layout (baselines to \SI{130}{\metre},
Section~\ref{sec:sweep}) increases the image DR by approximately
$20~\si{dB}$ because the more uniform 2D UV coverage suppresses PSF sidelobes
and eliminates the elongation artefact. This is the quantitative motivation for the
array upgrade discussed in the Conclusion.

\subsection{Interferometric Advantage for Burst Localisation}

A total-power (single-dish) radiometer can detect a burst but cannot localise it
on the solar disk. The compact interferometer adds a spatial dimension: the
fringe amplitude on the \SI{24}{\metre} baseline is sensitive to structures larger than
$\lambda/b\approx1.2^\circ$, larger than the solar diameter ($0.53^\circ$), meaning the
compact array detects but cannot resolve solar disk structure.
Although the compact east-west array cannot image the full disk, it can
distinguish between a burst originating at the disk centre (maximum fringe
amplitude) and one at the limb (reduced amplitude due to partial resolution). The
extended graded cross, reaching \SI{130}{\metre}, reduces the resolution element
to $\lambda/b_{\rm max}\approx0.22^\circ$, sufficient to localise a burst to a
fraction of the solar radius. Table~\ref{tab:detectability} and the dynamic-range
analysis together define the sensitivity envelope within which such localisation is
possible, and reinforce that the graded-cross layout of Section~\ref{sec:sweep} is
the one that best supports future burst-localisation science.

% ---------------------------------------------------------------------
\section{Antenna-Beam Fidelity and Its Effect on the Transit Power Profile}
\label{sec:beam-fidelity}

The sensitivity results of Section~\ref{sec:snr-analysis} assumed a specific primary
beam model. This section asks whether that choice matters: for a simulation that is
to be compared with real field measurements, the primary beam model must be accurate
at large zenith angles, not just at boresight. A common simplification, used by
several simplified aperture-synthesis simulators including \textsc{APSYNSIM}
(discussed in Section~\ref{sec:apsynsim}), is the Rayleigh-criterion Gaussian
approximation,
\begin{equation}
  B(\theta) = \exp\!\left(-\frac{\theta^2}{2\sigma_B^2}\right),
  \quad
  \sigma_B = \frac{1}{2.3548}\cdot\frac{1.22\,\lambda}{D_{\rm ap}},
  \label{eq:gauss-beam}
\end{equation}
where $D_{\rm ap}$ is the antenna aperture diameter (not to be confused with the
directivity $D$ of Eq.~\eqref{eq:directivity} in Chapter~\ref{ch:theory}) and
$\lambda$ is the wavelength. This captures the
main-lobe width but ignores all sidelobe structure, impedance-mismatch effects, and
the frequency-dependent evolution of the beam. Three specific failures of
Eq.~\eqref{eq:gauss-beam}, evaluated against the real 4NEC2 Yagi pattern of
Chapter~\ref{ch:antenna}, are relevant to this thesis.

\paragraph{Incorrect roll-off.}
The Gaussian falls to zero much faster than the real Yagi pattern at angles beyond
the main lobe. As shown in Table~\ref{tab:beam-comparison}, at
$\theta_z=60^\circ$ the Gaussian predicts only $0.063$ of the boresight response,
while the real 4NEC2 pattern retains $0.180$-a factor of $2.9\times$ higher.
Integrated over the 12-hour spring-equinox transit this discrepancy causes the
Gaussian to underestimate the total accumulated signal energy by approximately
$25\,\%$ ($\sim\!1~\si{dB}$), which would appear as a systematic flux-calibration
error in any power comparison with the field hardware.

\paragraph{Missing sidelobe contribution.}
A well-optimised Yagi has a first sidelobe approximately $13~\si{dB}$ below the
main lobe. When the Sun enters the first sidelobe (near the horizon at transit
edges), the Gaussian predicts essentially zero response. The 4NEC2 pattern retains
a non-zero response that would appear as a ghost fringe at a slightly different phase
than the main-lobe fringe-a signature that the Gaussian model would miss entirely.

\paragraph{No backlobe.}
The 4NEC2 simulation shows a back-hemisphere response of approximately
$-25~\si{dB}$ relative to the main lobe. While this is too faint to affect
calibration materially in a ground-mounted configuration, it becomes relevant if the
element is tilted or if a ground-plane reflection is to be included.

\begin{table}[H]
  \centering
  \caption[Gaussian versus 4NEC2 beam response at representative solar zenith angles]{Gaussian versus 4NEC2 beam response at representative solar zenith angles.\\[4pt]
  The normalised power response $A_n(\theta_z)$ is compared between the
  Rayleigh-criterion Gaussian approximation (Eq.~\eqref{eq:gauss-beam}) and
  the real 4NEC2 Yagi pattern at six zenith angles spanning the 12-hour
  transit. Both models share the same $52^\circ$ HPBW and agree to within
  3\% near transit; the ratio grows to $2.86\times$ at the $60^\circ$ horizon
  limit, where the Gaussian substantially underestimates the true beam
  response, leading to a systematic $\approx25\%$ ($\approx1$~dB) underestimate
  of integrated transit power.}
  \label{tab:beam-comparison}
  \renewcommand{\arraystretch}{1.3}
  \begin{tabular}{c c c c c}
    \toprule
    $\boldsymbol{\theta_z}$ & \textbf{Hour angle} & \textbf{Gaussian} $A_n$
      & \textbf{4NEC2} $A_n$ & \textbf{Ratio (4NEC2/Gauss)} \\
    \midrule
    $0^\circ$  & $0^\circ$ (transit) & 1.000 & 1.000 & 1.00 \\
    $10^\circ$ & $\pm15^\circ$       & 0.938 & 0.965 & 1.03 \\
    $20^\circ$ & $\pm30^\circ$       & 0.761 & 0.810 & 1.06 \\
    $30^\circ$ & $\pm45^\circ$       & 0.532 & 0.620 & 1.17 \\
    $45^\circ$ & $\pm64^\circ$       & 0.243 & 0.390 & 1.60 \\
    $60^\circ$ & $\pm84^\circ$       & 0.063 & 0.180 & 2.86 \\
    \bottomrule
  \end{tabular}
\end{table}

The most directly quantifiable impact is on the transit-averaged signal power
$\overline{P}_{\rm sim}$ (Eq.~\eqref{eq:pavg-sim}, used in Section~\ref{sec:real-comparison}).
The real beam predicts approximately $25\,\%$ more integrated energy over the
12-hour transit than the Gaussian approximation at the same peak gain. This
$\sim\!1~\si{dB}$ difference is the primary reason why the real 4NEC2 pattern, rather
than Eq.~\eqref{eq:gauss-beam}, was used throughout this thesis. Computationally, the
real beam adds negligible overhead: the \texttt{RegularGridInterpolator} evaluates the
pre-loaded $({\theta},{\phi})$ grid at any query point in sub-millisecond time, so there
is no performance argument for retaining the Gaussian approximation.

% ---------------------------------------------------------------------
\section{Comparison with Educational Simulators: \textsc{APSYNSIM}}
\label{sec:apsynsim}

The preceding sections validated the pipeline and used it for design and sensitivity
studies. Before turning to the hardware comparison, it is useful to place the
simulation environment in context against existing open-source aperture-synthesis
tools, so that the contribution of this work is clear. The most widely used such tool
is \textsc{APSYNSIM} (Aperture Synthesis Simulator), developed by
Mart\'i-Vidal at the Onsala Space Observatory for the Nordic ALMA Regional Centre
Interferometry School~\cite{apsynsim}.

\subsection{Overview and Feature Comparison}

\textsc{APSYNSIM} is a Python/Tkinter graphical application in which the user places
antennas interactively on a canvas and observes, in real time, how the UV-plane
coverage, the dirty beam, and the dirty image respond as elements are added, removed,
or repositioned. The tool ships with a library of well-known array configurations
(VLA, ALMA, WSRT, Golay, Golomb) and source templates (point sources, Gaussians,
uniform disks) and includes a H\"ogbom \textsc{clean}~\citep{hogbom1974} implementation for basic
deconvolution demonstration. As a classroom and workshop tool it is highly effective:
it builds the intuition that connects array layout to image quality without requiring
any knowledge of signal processing or calibration pipelines. Its design is,
however, deliberately simplified at the physics level: the primary beam is the
Rayleigh-criterion Gaussian of Eq.~\eqref{eq:gauss-beam} (Section~\ref{sec:beam-fidelity}),
which ignores sidelobe structure and impedance-mismatch effects, and the noise model
is an empirical Gaussian random vector scaled by a sensitivity parameter with no
reference to system temperature, bandwidth, or the Nyquist-Johnson relation. Real
observational data cannot be ingested; \textsc{APSYNSIM} operates exclusively on
synthetic source models.

Table~\ref{tab:apsynsim-comparison} summarises the resulting feature comparison.

\begin{table}[H]
  \centering
  \caption[Feature comparison: \textsc{APSYNSIM} versus this work's simulation pipeline]{Feature comparison: \textsc{APSYNSIM} versus this work's simulation pipeline.\\[4pt]
  The comparison covers interface, antenna beam model, noise treatment, data
  ingestion, coordinate geometry, hardware quantisation, batch processing,
  field validation, deconvolution, and reproducibility.}
  \label{tab:apsynsim-comparison}
  \renewcommand{\arraystretch}{1.35}
  \begin{tabularx}{\textwidth}{l X X}
    \toprule
    \textbf{Aspect} & \textbf{\textsc{APSYNSIM}} & \textbf{This Work} \\
    \midrule
    Interface
      & GUI only (Tkinter); not scriptable; not batch-processable
      & Jupyter notebook; fully scriptable; headless execution supported \\
    Antenna beam model
      & Rayleigh-criterion Gaussian (Eq.~\eqref{eq:gauss-beam})
      & Real 4NEC2 Yagi pattern from \texttt{.npz};
        bilinear \texttt{RegularGridInterpolator} \\
    Noise model
      & Empirical Gaussian: $\sigma_n \propto N_{\rm bas}^{-1/2}$; no physical basis
      & Johnson-Nyquist~\citep{nyquist1928}: $P_n = k_B T_{\rm sys}\Delta\nu$; thermodynamically grounded \\
    System temperature
      & Not modelled
      & $T_{\rm sys} = \SI{150}{\kelvin}$ (configurable per deployment) \\
    Receiver chain
      & None
      & Full Friis chain: $P_{\rm sig}=\tfrac{1}{2}A_e A_n(\theta_z) S_\nu \Delta\nu$ \\
    Data ingestion
      & Synthetic sky models only (\texttt{.model} parameter files)
      & Real archived SWS Learmonth \SI{610}{\mega\hertz} solar flux \\
    Coordinate geometry
      & Implicit; ENU positions entered directly
      & Explicit ENU $\rightarrow$ XYZ $\rightarrow$ $\uvw$~\citep{tms2017} \\
    Hardware quantisation
      & None
      & 16-bit IQ hex; LSB $\geq 3$ noise-floor validation for FPGA \\
    Batch processing
      & Not supported (interactive only)
      & 5-layout automated sweep; \texttt{.npz} + PNG export per geometry \\
    Field validation
      & None
      & Cross-validated vs 2-element analogue array at Xuan~Mai \\
    Deconvolution
      & H\"ogbom \textsc{clean} (interactive mask editing)
      & Not the focus; dirty image used as the analytic benchmark \\
    Reproducibility
      & Manual (session-dependent)
      & Full: \texttt{.ipynb} re-run reproduces every figure deterministically \\
    \bottomrule
  \end{tabularx}
\end{table}

\subsection{Noise Model: Physical Derivation versus Ad-Hoc Scaling}

The difference in noise treatment is central to the comparison. In
\textsc{APSYNSIM}, the noise added to each visibility is a free parameter
that the user scales until the image looks plausible; doubling $T_{\rm sys}$ or
halving the bandwidth has no effect because neither quantity exists in the model.
In the present pipeline the noise power derives from first principles: for a
receiver at system temperature $T_{\rm sys}$ terminated into resistance $R_{\rm load}$
over bandwidth $\Delta\nu$, the RMS magnitude of the complex Johnson-Nyquist~\citep{nyquist1928}
noise voltage at the load is
\begin{equation}
  V_{\rm noise} = \sqrt{k_B\,T_{\rm sys}\,\Delta\nu\,R_{\rm load}} = \sqrt{2}\,\sigma_n,
  \label{eq:johnson-rms}
\end{equation}
where $\sigma_n$ is the per-quadrature noise standard deviation already defined by
Eq.~\eqref{eq:noise-sigma} of Chapter~\ref{ch:pipeline}. For the simulation parameters
($T_{\rm sys}=\SI{150}{\kelvin}$, $\Delta\nu=\SI{1}{\mega\hertz}$,
$R_{\rm load}=\SI{75}{\ohm}$), $V_{\rm noise}\approx\SI{0.39}{\micro\volt}$ - a
thermodynamic lower bound, not a tunable parameter. Because the signal voltage
$V_{\rm rms}$ (Eq.~\eqref{eq:vrms} of Chapter~\ref{ch:theory}) is also computed from
the Friis chain using the real flux, the ratio $V_{\rm rms}/V_{\rm noise}$ is a
physically meaningful per-sample signal-to-noise ratio, a quantity that
\textsc{APSYNSIM} cannot produce.

\subsection{Workflow and Reproducibility}

\textsc{APSYNSIM} is designed for a single operator running one configuration
interactively: a new array geometry requires manually repositioning elements on the
GUI canvas, reselecting a source model, and reading beam and image quality off the
screen, with no guarantee that two sessions used identical settings. This is
appropriate for teaching but constrains scientific use, where exploring many dates,
baselines, and array geometries - potentially dozens of configurations in a single
study - is routine. In the present pipeline, every parameter is declared once in a
single configuration cell (Chapter~\ref{ch:pipeline}), and changing the observation
date or array layout requires editing only that cell before re-running the notebook
top to bottom. The parametric sweep of Section~\ref{sec:sweep} demonstrates the
result directly: five array geometries were processed automatically, with $uv$
coverage, dirty beams, and dirty images exported without manual intervention.
Table~\ref{tab:workflow-compare} summarises the resulting difference in workflow for a
representative multi-geometry, multi-date design study.

\begin{table}[H]
  \centering
  \caption[Workflow comparison: \textsc{APSYNSIM} versus this work's pipeline]{Workflow comparison: \textsc{APSYNSIM} versus this work's pipeline.\\[4pt]
  A typical instrument-design study covering five array geometries on two
  observation dates. Steps include specifying the array, loading source data,
  setting observation parameters, running the simulation, exporting results,
  reproducing results later, and changing the observation date.}
  \label{tab:workflow-compare}
  \renewcommand{\arraystretch}{1.3}
  \begin{tabularx}{\textwidth}{l X X}
    \toprule
    \textbf{Step} & \textbf{\textsc{APSYNSIM}} & \textbf{This Work} \\
    \midrule
    Specify array
      & Click antenna positions on GUI canvas (manual, session-by-session)
      & Edit one list of ENU coordinates in Cell~0; or call a preset function \\
    Load source data
      & Select from built-in menu; adjust Gaussian parameters interactively
      & Automatic: \texttt{sws.load\_day(NGAY)} reads real flux from file \\
    Set observation parameters
      & Adjust GUI sliders for latitude, HA, declination, wavelength
      & Edit 2-3 variables in Cell~0; all cells inherit automatically \\
    Run simulation
      & Immediate (GUI redraws on each mouse event)
      & Run all cells: $<\SI{2}{\minute}$ for 50\,000 samples \\
    Export results
      & Screenshot / manual FITS save (APSYNTRUE only)
      & Automatic: \texttt{.npz} + PNG per geometry, summary \texttt{.txt} \\
    Reproduce later
      & Re-run manually; no guaranteed parameter record
      & Re-execute \texttt{.ipynb}; identical results guaranteed \\
    Change observation date
      & Full manual re-setup required
      & Edit \texttt{NGAY} and \texttt{DEC\_SUN\_DEG} in Cell~0; re-run \\
    \bottomrule
  \end{tabularx}
\end{table}

Because all parameters live in one cell, a parameter sensitivity study is
similarly straightforward: changing $T_{\rm sys}$ from \SI{150}{\kelvin} to
\SI{200}{\kelvin} reduces the per-dump SNR by $\approx1.2~\si{dB}$, and changing
\texttt{LAT\_DEG} to $-26.7^\circ$ would simulate the Murchison Radio-astronomy
Observatory with the ENU$\rightarrow$XYZ rotation and transit window updating
automatically. None of these changes requires modifying more than one or two lines of
code, and none are possible in a GUI-only tool without full manual re-interaction.
This scriptability, rather than any single figure of merit, is what separates the
present pipeline from an educational visualisation tool and makes it usable as a
design instrument for the studies in Sections~\ref{sec:sweep}-\ref{sec:beam-fidelity}.
A companion browser-based interface that packages these same outputs into an
accessible, zero-installation viewer, without introducing any new physical model, is
described in the Appendix (Project Repository and Simulation Guide).

% ---------------------------------------------------------------------
\section{Limitations of the Current Implementation}
\label{sec:limitations}

The validation, design, and sensitivity results above establish confidence in the
pipeline, but that confidence is bounded by a specific set of simplifications. For
completeness, the principal simplifications and their expected impact on results
are listed below; each is revisited as a concrete extension in the Future Work
section of the Conclusion.

\begin{enumerate}[leftmargin=1.8em]
  \item \textbf{Single spectral channel.} The simulation models a single
        $\SI{1}{\mega\hertz}$ channel at $\SI{610}{\mega\hertz}$. Bandwidth
        smearing-which causes decorrelation at angular offsets from the phase
        centre proportional to $\Delta\nu/\nu$-is not modelled. For
        $\Delta\nu/\nu = 10^6/610\times10^6 \approx 1.6\times10^{-3}$ and
        baselines below $\SI{130}{\metre}$, the smearing at a $0.5^\circ$ offset
        is $\lesssim2\,\%$ of the fringe amplitude, which is negligible for the
        present validation study.
  \item \textbf{No ionospheric phase screen.} Differential ionospheric
        phase between antennas scales as $\Delta\phi \approx \Delta{\rm TEC}/\nu^2$,
        where $\Delta{\rm TEC}$ is the differential total electron content. For
        baselines below $\SI{130}{\metre}$ at $\SI{610}{\mega\hertz}$, the expected
        differential phase is $\lesssim1^\circ$ during undisturbed conditions, which
        is below the current precision requirement.
  \item \textbf{Independent-element assumption.} The mutual-coupling study of
        Chapter~\ref{ch:antenna} showed coupling below $-78~\si{dB}$ beyond
        $\SI{8}{\metre}$. No Jones-matrix correction is applied even at the shortest
        $\SI{8}{\metre}$ baseline, where the coupling could in principle contribute
        a $\sim10^{-4}$ fractional amplitude error-far below any other error source.
  \item \textbf{No deconvolution.} The imaging chain produces a dirty image.
        \textsc{clean} or a compressed-sensing solver would be required to recover
        the true sky brightness and to achieve noise-limited dynamic range.
  \item \textbf{Narrow-field assumption.} The DFT imaging assumes a flat sky
        (small $w$-term). The solar diameter ($0.53^\circ$) is well within the
        regime where the $w$-approximation is valid for baselines below
        $\SI{130}{\metre}$, but extension to wider fields or longer baselines would
        require $w$-projection.
  \item \textbf{Perfect clock synchronisation.}
        The voltage model of Eq.~\eqref{eq:vk} applies a common, deterministic
        time axis to all antennas (\texttt{CLOCK\_OFFSET\_NS}$=[0,0,0,0]$~ns).
        In real hardware each digitiser runs on its own oscillator; a differential
        clock drift $\xi$ between any two antennas adds a time-varying
        phase to their shared visibility and can reduce the time-averaged fringe
        amplitude to zero within a single integration dump-independently of SNR.
        A quantitative sensitivity study is presented in
        Section~\ref{sec:clock-offset}.
\end{enumerate}

% ---------------------------------------------------------------------
\section{Diagnostic Comparison with Field Measurement}
\label{sec:real-comparison}

The preceding sections verify the simulation against \emph{analytic} predictions
(fringe period, $\jinc$ visibility, PSF identity) and use it for design and
sensitivity studies. This section asks a different, harder question: how does the
simulation compare with an actual hardware recording? It extends the comparison
to an \emph{experimental} recording from the two-element analogue interferometer
deployed at Hoa~Lac on 20~March~2025-the same date used throughout this chapter.
What is compared is deliberately narrow: the gain-independent solar excess above the
noise floor, not the raw voltage waveform, a phase-coherent visibility, or a
calibrated image, since the hardware gain chain was not independently calibrated end
to end. This comparison is therefore a \emph{hardware sanity check}, not a full
end-to-end validation.

\subsection{Digital versus Analogue Interferometry}

In \emph{analogue interferometry} the two antenna signals are combined at radio
frequency by a coaxial power combiner. The combined signal drives a broadband power
detector; the output is a slowly varying DC voltage proportional to the total received
power including the fringe modulation. The phase relationship between the two signals
is fixed by the physical cable lengths and the geometry of the array; it cannot be
adjusted post-detection. The instrument is a correlating radiometer realised entirely
in hardware.

In \emph{digital interferometry} (the approach of this thesis) the two signals are
digitised independently as complex IQ baseband voltages. The cross-correlation is
computed in software, allowing arbitrary post-processing: fringe stopping, phase
calibration, baseline selection, and imaging. The key advantage is that a single
recording can be reprocessed with different array parameters, different imaging
algorithms, or additional baselines added later, without any physical change to the
hardware.

\subsection{Average Power on 20~March~2025}

The simulation computes the beam-weighted transit-average signal power over the
observation window $|H|\leq30^\circ$ (approximately 10:00-14:00 local time,
matching the hardware recording window):
\begin{equation}
  \overline{P}_{\rm sim} = \frac{1}{N}\sum_{n=0}^{N-1}
    S_\nu(t_n)\,A_e\,A_n\!\bigl(\theta_z(t_n)\bigr)\,\Delta\nu,
  \label{eq:pavg-sim}
\end{equation}
where $A_n(\theta_z)$ is the normalised power beam pattern from the 4NEC2 simulation
(Section~\ref{sec:beam-fidelity}). Because the absolute power at the hardware receiver
port depends on the unknown total gain of the analogue chain (LNA, cables,
downconverter), a direct watt-to-watt comparison is not feasible without an
end-to-end gain calibration. Instead, the comparison uses the \emph{solar excess above
the thermal noise floor},
\begin{equation}
  \mathcal{E} = 10\log_{10}\!\left(1 + \frac{\overline{P}_{\rm sim}}{P_n}\right)
  \quad[\mathrm{dB}],
  \label{eq:excess}
\end{equation}
where $P_n$ is the Johnson-Nyquist noise power of Eq.~\eqref{eq:johnson} in
Chapter~\ref{ch:theory}. This quantity is gain-independent: any common
multiplicative factor in the hardware chain cancels between signal and noise,
leaving only the ratio $\overline{P}_{\rm sim}/P_n$.
From the hardware trace recorded at Hoa~Lac on 20~March~2025, the corresponding
measurement is read directly as the difference between the transit envelope peak and
the pre-transit noise baseline (both in the same dB scale of the recording).

Table~\ref{tab:power-comparison} compares $\mathcal{E}$ from simulation and measurement.

\begin{table}[H]
  \centering
  \caption[Solar excess for 20~March~2025: simulation versus hardware measurement]{Solar excess for 20~March~2025: simulation versus hardware measurement.\\[4pt]
  The solar excess $\mathcal{E}$ (Eq.~\eqref{eq:excess}) is the ratio of
  time-averaged signal power to the thermal noise floor, expressed in dB.
  Simulation: 4NEC2 beam model, archived SWS Learmonth flux,
  $T_{\rm sys}=\SI{150}{\kelvin}$, transit window $|H|\leq30^\circ$.
  Hardware: two-element Yagi analogue interferometer at Hoa~Lac on
  20~March~2025, operated without a dedicated low-noise amplifier.
  The factor-of-$23.8\times$ SNR gap implies an effective hardware system
  temperature of approximately \SI{3600}{\kelvin}.}
  \label{tab:power-comparison}
  \renewcommand{\arraystretch}{1.3}
  \begin{tabularx}{\textwidth}{l c >{\raggedright\arraybackslash}X}
    \toprule
    \textbf{Source} & \textbf{Solar excess $\mathcal{E}$ (dB)} & \textbf{Notes} \\
    \midrule
    Simulation (this work)
      & $0.586$
      & Eq.~\eqref{eq:excess}; $\overline{P}_{\rm sim}/P_n=14.3\,\%$ \\
    Hardware (Hoa~Lac, 20 Mar 2025)
      & $0.026$
      & Peak envelope minus baseline; $P_{\rm sig}/P_n=0.60\,\%$ \\
    Ratio (linear SNR)
      & $23.8\times$
      & Implies $T_{\rm sys,eff}\approx\SI{3600}{\kelvin}$ \\
    \bottomrule
  \end{tabularx}
\end{table}

\subsection{Interpretation}

Table~\ref{tab:power-comparison} reveals a factor of $23.8\times$ between the
simulated and measured solar excess. This discrepancy is not evidence against the
Friis model; it is best read as a measurement of the \emph{effective system
temperature gap} between the idealised pipeline ($T_{\rm sys}=\SI{150}{\kelvin}$) and
the real consumer-grade hardware. The real receiver is a software-defined radio (SDR)
tuner connected through passive coaxial splitters and cables, with no dedicated
low-noise amplifier (LNA), so its effective system temperature is set mainly by the
receiver noise figure $\mathrm{NF}$ and cable loss rather than by the antenna or sky
temperature. A typical RTL-SDR at \SI{610}{\mega\hertz} has
$\mathrm{NF}\approx8$--$12~\si{\decibel}$, corresponding to a receiver noise
temperature $T_{\rm rx}=(10^{\mathrm{NF}/10}-1)\,T_0\approx1\,500$--$4\,300~\si{\kelvin}$
for a reference temperature $T_0=\SI{290}{\kelvin}$; a few decibels of cable loss
raise this further, into the same thousands-of-kelvin range obtained by inverting the
measured ratio of Table~\ref{tab:power-comparison} directly:
\begin{equation}
  T_{\rm sys,eff} = T_{\rm sys,model}
    \times\frac{(\overline{P}_{\rm sim}/P_n)_{\rm sim}}
               {(\overline{P}_{\rm sig}/P_n)_{\rm meas}}
  = 150\times23.8 \approx \SI{3\,600}{\kelvin},
  \label{eq:tsys-backcalc}
\end{equation}
which is the effective system temperature used in the remainder of this section, and
is fully consistent with SDR hardware operating at UHF without an LNA and with
typical cable losses of $2$-$4~\si{\decibel}$.

This comparison is therefore consistent with the signal-level model used throughout
the simulation, and is corroborated independently by the fringe-period and
$\jinc$-visibility benchmarks of Sections~\ref{sec:fringe-results}
and~\ref{sec:vis-results}: it indicates that the field setup was receiver-noise
limited, rather than pointing to an error in the Friis chain itself. Replacing the
consumer SDR with an LNA of noise figure $\lesssim\SI{1}{\decibel}$
($T_{\rm rx}\approx\SI{75}{\kelvin}$) would reduce $T_{\rm sys,eff}$ to
$\approx\SI{300}{\kelvin}$ and increase the measured solar excess to
$\approx0.026\times(3\,600/300)\approx\SI{0.31}{\decibel}$, an improvement of
$12\times$ in SNR, comparable to the simulation model. Beyond total power, the
digital pipeline recovers information inaccessible to the analogue chain: the
$\sim$14-minute fringe modulation and the $\jinc$-shaped visibility amplitude versus
baseline length. This comparison - total power envelope (analogue, broad bell curve)
versus fringe pattern and visibility (digital) - is the clearest demonstration of the
complementary value of digital interferometry for solar burst science, provided the
receiver-noise and clock-coherence gaps identified here and in
Section~\ref{sec:clock-offset} are closed.

% ---------------------------------------------------------------------
\section{Clock-Coherence Requirement}
\label{sec:clock-offset}

Section~\ref{sec:real-comparison} attributed the $23.8\times$ SNR gap to receiver
system temperature alone. This section asks whether that is the only barrier: the
voltage model of Chapter~\ref{ch:pipeline} treats all antennas as sharing a perfect
common clock (\texttt{CLOCK\_OFFSET\_NS}$=[0,0,0,0]$~ns). Using the clock-offset
injection facility (Eq.~\eqref{eq:clock-phase}), this section quantifies how
realistic oscillator imperfections would degrade interferometric measurements, and
places the hardware coherence requirement on a quantitative footing.

\subsection{Static Offset: Phase Rotation Without Amplitude Loss}

A static clock offset $\tau_{ij}$ between the digitisers of antennas $i$ and $j$
introduces a constant phase error on the baseline visibility
(Eq.~\eqref{eq:vis-clock}):
\begin{equation}
  \Delta\phi_{ij} = 2\pi f_c\,\tau_{ij}
  \;\;\Longrightarrow\;\;
  \frac{\Delta\phi_{ij}}{1~\mathrm{ns}}
    = 2\pi\times\SI{610}{\mega\hertz}\times10^{-9}~\mathrm{s}
    \approx 219.6^\circ\,\mathrm{ns}^{-1}.
  \label{eq:static-phase}
\end{equation}
Table~\ref{tab:clock-static} evaluates Eq.~\eqref{eq:static-phase} for
representative offsets.

\begin{table}[H]
  \centering
  \caption[Static clock offset versus visibility phase error at 610~MHz]{Static clock offset versus visibility phase error at 610~MHz.\\[4pt]
  A static differential clock offset $\tau_{ij}$ between two digitisers
  introduces a constant phase rotation $\Delta\phi_{ij}=2\pi f_c\tau_{ij}$
  on the corresponding baseline visibility; the amplitude is unaffected.
  At $f_c=\SI{610}{\mega\hertz}$, even a 0.1~ns offset produces a $22^\circ$
  phase error, sufficient to displace a simulated point source by more than
  its own diameter along the 8-m baseline direction.}
  \label{tab:clock-static}
  \renewcommand{\arraystretch}{1.25}
  \begin{tabularx}{\textwidth}{c c >{\raggedright\arraybackslash}X}
    \toprule
    \textbf{Offset $\tau_{ij}$} & \textbf{Phase error $\Delta\phi_{ij}$}
      & \textbf{Dirty-image effect} \\
    \midrule
    $0$~ns    & $0^\circ$    & No effect; phase centre correct \\
    $0.1$~ns  & $22^\circ$   & Mild apparent source position shift \\
    $0.5$~ns  & $110^\circ$  & Source displaced along baseline direction \\
    $1$~ns    & $220^\circ$  & Large position error; image artefacts \\
    $1.64$~ns & $360^\circ$  & Full phase cycle; source wraps to original position \\
    \bottomrule
  \end{tabularx}
\end{table}

A static phase error shifts the apparent source position along the baseline
direction by $\Delta\theta\approx\Delta\phi_{ij}/(2\pi\,b_{ij}/\lambda)$. For the
\SI{8}{\metre} baseline ($b_{ij}/\lambda\approx16.3$) and a $1$~ns offset
($\Delta\phi_{ij}\approx3.83$~rad), this gives $\Delta\theta\approx2.1^\circ$-larger
than the solar diameter ($0.53^\circ$): a single-nanosecond timing error places the
simulated Sun at the wrong sky position by nearly four solar diameters. Static offsets
can be removed by a point-source phase calibration, but only if a suitable calibrator
transit is available.

\subsection{Clock Drift: Fringe Washing and the Coherence Time}

A more fundamental degradation arises when the two oscillators run at slightly
different \emph{rates}. Let $\xi$ (dimensionless, in parts per million, ppm)
be the fractional frequency offset between the digitisers of antennas $i$ and $j$ -
not to be confused with the stochastic source process $\varepsilon(t)$ of
Eq.~\eqref{eq:vk} in Chapter~\ref{ch:pipeline}, which is an unrelated quantity. The
differential delay then grows as $\tau_{ij}(t)=\tau_{ij}(0)+\xi\,t$, and the
instantaneous visibility phase advances continuously. Time-averaging over one
integration dump of duration $\Delta t_{\mathrm{dump}}=\SI{43.2}{\second}$ reduces the
fringe amplitude by the coherence factor
\begin{equation}
  \eta_{\mathrm{coh}}
  = \left|\operatorname{sinc}\!\bigl(f_c\,\xi\,\Delta t_{\mathrm{dump}}\bigr)\right|,
  \label{eq:coherence}
\end{equation}
where $\operatorname{sinc}(x)\equiv\sin(\pi x)/(\pi x)$, falling to zero (complete
fringe washing) when $f_c\,\xi\,\Delta t_{\mathrm{dump}}=1$, defining the
\emph{coherence time}
\begin{equation}
  \Delta t_{\mathrm{coh}} = \frac{1}{f_c\,\xi}.
  \label{eq:tcoh}
\end{equation}

Table~\ref{tab:clock-drift} evaluates Eqs.~\eqref{eq:coherence} and
\eqref{eq:tcoh} for oscillator accuracies representative of consumer and
laboratory hardware at $f_c=\SI{610}{\mega\hertz}$,
$\Delta t_{\mathrm{dump}}=\SI{43.2}{\second}$.

\begin{table}[H]
  \centering
  \caption[Clock drift versus fringe coherence factor at 610~MHz (43.2-s dump)]{Clock drift versus fringe coherence factor at 610~MHz (43.2-s dump).\\[4pt]
  The fringe coherence factor $\eta_{\rm coh}$ (Eq.~\eqref{eq:coherence})
  is evaluated as a function of oscillator fractional frequency offset
  $\xi$ for a 43.2-second integration dump at $f_c=\SI{610}{\mega\hertz}$.
  Consumer SDR crystals (1~ppm drift) produce complete fringe washing
  ($\eta_{\rm coh}\approx0$) within a single dump; maintaining
  $\eta_{\rm coh}>0.9$ requires a GPSDO or equivalent precision oscillator
  with $\xi\lesssim10$~ppt.}
  \label{tab:clock-drift}
  \renewcommand{\arraystretch}{1.25}
  \begin{tabularx}{\textwidth}{l c c c >{\raggedright\arraybackslash}X}
    \toprule
    \textbf{Drift $\xi$} & \textbf{$\Delta t_{\rm coh}$}
      & \textbf{$f_c\xi\Delta t_{\rm dump}$}
      & \textbf{$\eta_{\rm coh}$} & \textbf{Oscillator class} \\
    \midrule
    $1$~ppm         & $1.6$~ms         & $26\,352$  & $\approx 0$   & Consumer SDR crystal \\
    $0.1$~ppm       & $16$~ms          & $2\,635$   & $\approx 0$   & TCXO (temperature compensated) \\
    $1$~ppb         & $1.6$~s          & $26.4$     & $\approx 0$   & OCXO (oven controlled) \\
    $10$~ppt        & $164$~s          & $0.264$    & $0.88$        & Precision OCXO \\
    $1$~ppt         & $1\,640$~s       & $0.026$    & $\approx 1$   & GPSDO (GPS-disciplined) \\
    \bottomrule
  \end{tabularx}
\end{table}

The requirement for $\eta_{\mathrm{coh}}>0.9$ over the \SI{43.2}{\second} dump is
$\xi<10$~ppt ($10^{-11}$), achievable with a GPS-disciplined oscillator
(GPSDO) or a precision oven-controlled crystal (OCXO) with external \SI{10}{\mega\hertz}
reference injection. Consumer oscillators-including the crystal in a typical
RTL-SDR, rated at approximately $0.5$-$1$~ppm-produce
$\Delta t_{\mathrm{coh}}\lesssim\SI{3}{\milli\second}$, more than four orders of
magnitude shorter than $\Delta t_{\mathrm{dump}}$: a consumer SDR therefore produces
zero fringe amplitude within the very first integration dump, regardless of SNR.

\subsection{Implication for the Hardware Comparison}

The clock-offset analysis adds a second, \emph{independent} reason why digital
cross-correlation with unsynchronised SDRs cannot recover a fringe, alongside the
system-temperature gap of Section~\ref{sec:real-comparison}: even if
$T_{\mathrm{sys,eff}}$ were reduced to the idealised \SI{150}{\kelvin} by adding a
dedicated LNA, the time-averaged fringe amplitude would remain zero because
$\eta_{\mathrm{coh}}\approx0$ for any consumer oscillator. These two effects are
therefore complementary barriers, and the minimum hardware specification for digital
solar interferometry at \SI{610}{\mega\hertz} with \SI{43.2}{\second} dumps requires
addressing \emph{both} simultaneously:
\begin{enumerate}[leftmargin=1.8em]
  \item a low-noise amplifier with noise figure $\lesssim\SI{1}{\decibel}$
        ($T_{\mathrm{rx}}\approx\SI{75}{\kelvin}$), reducing
        $T_{\mathrm{sys,eff}}$ from $\approx\SI{3\,600}{\kelvin}$ to
        $\approx\SI{300}{\kelvin}$; and
  \item a common clock reference, a GPSDO or a \SI{10}{\mega\hertz}
        distribution amplifier, injected into all digitiser clock inputs,
        bringing $\xi$ below $10$~ppt and
        $\eta_{\mathrm{coh}}>0.9$.
\end{enumerate}
Meeting both conditions would increase the measured solar excess from the
hardware-measured $\mathcal{E}_{\mathrm{meas}}=0.026$~dB to
$\approx0.026\times(3\,600/300)\times\eta_{\mathrm{coh}}
\approx\SI{0.31}{\decibel}$, consistent with the simulation prediction of
$\mathcal{E}_{\mathrm{sim}}=0.586$~dB.

% ---------------------------------------------------------------------
\section{Chapter Conclusion}

The results of this chapter show that the simulation pipeline is useful not only as a
numerical demonstration, but as a design tool for a future low-cost solar radio
interferometer, following the validation-to-deployment logic set out in
Section~\ref{sec:result-strategy}.

First, the analytic validation tests (Sections~\ref{sec:fringe-results}-\ref{sec:dirty-results})
confirm that the main signal-processing blocks behave as expected: the simulated
fringe period verifies the geometric-delay and $\uvw$-coordinate implementation, the
visibility-amplitude trend verifies that the resolved solar disk is correctly
represented, and the dirty-image comparison verifies that the imaging stage produces
the expected sky-convolved-with-PSF result. Together these validate the internal
consistency of the pipeline and justify using it for the design work that follows.

Second, the array-geometry study (Section~\ref{sec:sweep}) shows why simulation is
necessary before field deployment. A compact east-west layout can demonstrate
fringes but gives limited two-dimensional localisation; the graded-cross and
extended layouts provide better $uv$-coverage and a more useful synthesised beam, but
also demand more careful calibration and, as Section~\ref{sec:clock-offset} shows,
clock coherence. The DLITE-inspired triangle-plus-centre alternative (Case~1.75) shows
that non-redundant $uv$-sampling, not baseline length, is what lets exactly four
elements recover useful two-dimensional resolution - the same minimum-redundancy
principle that motivates a future four-element LWA-style deployment. The best antenna
layout is therefore not determined by baseline length alone, but by the combined
behaviour of $uv$-coverage, PSF shape, sidelobe level, and practical deployability - a
conclusion reinforced by the sensitivity and dynamic-range analysis of
Section~\ref{sec:snr-analysis} and by the beam-fidelity result of
Section~\ref{sec:beam-fidelity}, where replacing the Gaussian primary beam with the
real 4NEC2 pattern changes the predicted transit power by $\approx25\,\%$.

Third, the comparative and diagnostic results identify the concrete barriers between
simulation and measurement. Section~\ref{sec:apsynsim} shows that the pipeline goes
beyond educational visualisation by combining physically grounded noise, real data
ingestion, real antenna beams, hardware export, and full reproducibility - the
capabilities needed to run the design studies above at all. Section~\ref{sec:real-comparison}
compares the simulated and measured solar excess for 20~March~2025
($\mathcal{E}_{\rm sim}=0.586$~dB versus $\mathcal{E}_{\rm meas}=0.026$~dB, a
$23.8\times$ gap in linear SNR) and finds it consistent with an elevated effective
system temperature in the consumer SDR receiver chain
($T_{\rm sys,eff}\approx\SI{3\,600}{\kelvin}$); this diagnostic comparison should not
be read as a full end-to-end hardware validation, which would require an
independently gain-calibrated measurement. Section~\ref{sec:clock-offset} adds a
second, independent barrier: a consumer crystal oscillator with
$\xi\approx1$~ppm gives a coherence time of only $\approx\SI{1.6}{\milli\second}$,
four orders of magnitude shorter than the \SI{43.2}{\second} dump, so unsynchronised
digitisers yield zero fringe amplitude regardless of SNR.

The practical outcome of this chapter is therefore a set of deployment
recommendations rather than a single number: use a calibrated low-noise receiver
chain ($\mathrm{NF}\lesssim\SI{1}{\decibel}$), distribute a common clock reference
($\xi<10$~ppt, e.g.\ a GPSDO), preserve the real antenna primary-beam model
rather than a Gaussian approximation, and choose the array layout based on
$uv$-coverage and sidelobe behaviour rather than on baseline length alone. These
recommendations, together with the six limitations of Section~\ref{sec:limitations},
define the concrete next steps for turning the present detection instrument into a
working digital solar imaging interferometer.
\clearpage
